% GENERATED FILE. Edit canonical lessons, then run npm run generate:books.
% canonical-source-hash: fnv1a64:06ac9bc29dc8e356
% canonical-lessons: TE-C200-podupu-ceyu, TE-C200-encuko, TE-C200-prayatnincu, TE-C200-vaddincu, TE-C200-namalu

\chapter{To save, To choose, To try, To serve, To chew}
\label{ch:te-to-save-to-choose-to-try-to-serve-to-chew}
\hlchaptermodality{200}

\begin{chapteropening}
\textbf{By the end of this chapter:} \emph{I can say to save, to choose, to try, to serve and to chew.}

Complete the last lesson of chapter 200: I can say to save, to choose, to try, to serve and to chew.
\end{chapteropening}

\section[\texorpdfstring{\te{పొదుపు} \te{చేయు} (podupu cēyu)}{పొదుపు చేయు (podupu cēyu)}]{\te{పొదుపు} \te{చేయు} (podupu cēyu) — to save (money)}

\label{lesson:TE-C200-podupu-ceyu}



\begin{quote}
Before the new one: say the Telugu for to fold, then the Telugu for to arrange.
\end{quote}

\subsection*{You'll want to know: \te{పొదుపు} \te{చేయు}}
\textbf{\te{పొదుపు} \te{చేయు}} — \emph{podupu cēyu} — \textquotedblleft{}to save (money)\textquotedblright{}.

\begin{grammarlens}[title={What you've built}]
One more word for this chapter.
\end{grammarlens}

\subsection*{Your turn}
\begin{itemize}
\raggedright
  \item \emph{Say it:} \emph{podupu cēyu}
  \item \emph{Say it:} \emph{podupu cēyu}, once more
  \item \emph{Say it:} say \emph{podupu cēyu}
  \item \emph{Recall:} say the Telugu for to fold, then the Telugu for to arrange, then say \emph{podupu cēyu} again
\end{itemize}

\begin{tcolorbox}[breakable,colback=teal!4,colframe=teal!35!black,title={Before you move on}]
What does \te{పొదుపు} \te{చేయు} mean? (\textquotedblleft{}To save (money)\textquotedblright{}.) Say it once more.
\end{tcolorbox}
\section[\texorpdfstring{\te{ఎంచుకో} (eñcukō)}{ఎంచుకో (eñcukō)}]{\te{ఎంచుకో} (eñcukō) — to choose}

\label{lesson:TE-C200-encuko}



\begin{quote}
Before the new one: say the Telugu for to arrange, then the Telugu for to save.
\end{quote}

\subsection*{You'll want to know: \te{ఎంచుకో}}
\textbf{\te{ఎంచుకో}} — \emph{eñcukō} — \textquotedblleft{}to choose\textquotedblright{}.

\begin{grammarlens}[title={What you've built}]
One more word for this chapter.
\end{grammarlens}

\subsection*{Your turn}
\begin{itemize}
\raggedright
  \item \emph{Say it:} \emph{eñcukō}
  \item \emph{Say it:} \emph{eñcukō}, once more
  \item \emph{Say it:} say \emph{eñcukō}
  \item \emph{Recall:} say the Telugu for to arrange, then the Telugu for to save, then say \emph{eñcukō} again
\end{itemize}

\begin{tcolorbox}[breakable,colback=teal!4,colframe=teal!35!black,title={Before you move on}]
What does \te{ఎంచుకో} mean? (\textquotedblleft{}To choose\textquotedblright{}.) Say it once more.
\end{tcolorbox}
\section[\texorpdfstring{\te{ప్రయత్నించు} (prayatniñcu)}{ప్రయత్నించు (prayatniñcu)}]{\te{ప్రయత్నించు} (prayatniñcu) — to try}

\label{lesson:TE-C200-prayatnincu}



\begin{quote}
Before the new one: say the Telugu for to save, then the Telugu for to choose.
\end{quote}

\subsection*{You'll want to know: \te{ప్రయత్నించు}}
\textbf{\te{ప్రయత్నించు}} — \emph{prayatniñcu} — \textquotedblleft{}to try\textquotedblright{}.

\begin{grammarlens}[title={What you've built}]
One more word for this chapter.
\end{grammarlens}

\subsection*{Your turn}
\begin{itemize}
\raggedright
  \item \emph{Say it:} \emph{prayatniñcu}
  \item \emph{Say it:} \emph{prayatniñcu}, once more
  \item \emph{Say it:} say \emph{prayatniñcu}
  \item \emph{Recall:} say the Telugu for to save, then the Telugu for to choose, then say \emph{prayatniñcu} again
\end{itemize}

\begin{tcolorbox}[breakable,colback=teal!4,colframe=teal!35!black,title={Before you move on}]
What does \te{ప్రయత్నించు} mean? (\textquotedblleft{}To try\textquotedblright{}.) Say it once more.
\end{tcolorbox}
\section[\texorpdfstring{\te{వడ్డించు} (vaḍḍiñcu)}{వడ్డించు (vaḍḍiñcu)}]{\te{వడ్డించు} (vaḍḍiñcu) — to serve (food)}

\label{lesson:TE-C200-vaddincu}



\begin{quote}
Before the new one: say the Telugu for to choose, then the Telugu for to try.
\end{quote}

\subsection*{You'll want to know: \te{వడ్డించు}}
\textbf{\te{వడ్డించు}} — \emph{vaḍḍiñcu} — \textquotedblleft{}to serve (food)\textquotedblright{}.

\begin{grammarlens}[title={What you've built}]
One more word for this chapter.
\end{grammarlens}

\subsection*{Your turn}
\begin{itemize}
\raggedright
  \item \emph{Say it:} \emph{vaḍḍiñcu}
  \item \emph{Say it:} \emph{vaḍḍiñcu}, once more
  \item \emph{Say it:} say \emph{vaḍḍiñcu}
  \item \emph{Recall:} say the Telugu for to choose, then the Telugu for to try, then say \emph{vaḍḍiñcu} again
\end{itemize}

\begin{tcolorbox}[breakable,colback=teal!4,colframe=teal!35!black,title={Before you move on}]
What does \te{వడ్డించు} mean? (\textquotedblleft{}To serve (food)\textquotedblright{}.) Say it once more.
\end{tcolorbox}
\section[\texorpdfstring{\te{నమలు} (namalu)}{నమలు (namalu)}]{\te{నమలు} (namalu) — to chew}

\label{lesson:TE-C200-namalu}



\begin{quote}
Before the new one: say the Telugu for to try, then the Telugu for to serve.
\end{quote}

\subsection*{You'll want to know: \te{నమలు}}
\textbf{\te{నమలు}} — \emph{namalu} — \textquotedblleft{}to chew\textquotedblright{}.

\begin{grammarlens}[title={What you've built}]
One more word for this chapter.
\end{grammarlens}

\subsection*{Your turn}
\begin{itemize}
\raggedright
  \item \emph{Say it:} \emph{namalu}
  \item \emph{Say it:} \emph{namalu}, once more
  \item \emph{Say it:} say \emph{namalu}
  \item \emph{Recall:} say the Telugu for to try, then the Telugu for to serve, then say \emph{namalu} again
\end{itemize}

\begin{tcolorbox}[breakable,colback=teal!4,colframe=teal!35!black,title={Before you move on}]
What does \te{నమలు} mean? (\textquotedblleft{}To chew\textquotedblright{}.) Say it once more.
\end{tcolorbox}
